\section{Temporal Siamese Network}

\label{sec:tsn}

\subsection{Architecture Overview}

The TSN processes a temporal sequence of multi-spectral satellite images to detect ecological changes.
Unlike bi-temporal methods that compare only two images, TSN ingests a sequence of $T$ images spanning the monitoring period, enabling robust change detection that distinguishes persistent land use change from transient events (cloud shadows, seasonal variation, short-term flooding).

The architecture consists of three components:

\begin{enumerate}[leftmargin=*]
    \item \textbf{Spatial Encoder:} A shared convolutional backbone that extracts spatial features from each image independently.
    \item \textbf{Temporal Attention Module:} A transformer-based module that models temporal dependencies across the image sequence, identifying persistent changes while suppressing transient artifacts.
    \item \textbf{Change Decoder:} A fully convolutional decoder that produces pixel-level change maps with class labels.
\end{enumerate}

\subsection{Spatial Encoder}

The spatial encoder $E_\psi$ processes each multi-spectral image $\mathbf{x}_t \in \mathbb{R}^{H \times W \times B}$ (where $B$ is the number of spectral bands) to produce a feature map:

\begin{equation}
    \mathbf{F}_t = E_\psi(\mathbf{x}_t) \in \mathbb{R}^{h \times w \times d}
\end{equation}

We use a ResNet-34 backbone pretrained on satellite imagery, modified to accept $B$-channel input (10 bands for Sentinel-2).
The encoder is shared across all timesteps (Siamese structure), ensuring consistent feature extraction.

\subsection{Temporal Attention Module}

The temporal attention module processes the sequence of feature maps $\{\mathbf{F}_1, \ldots, \mathbf{F}_T\}$ to identify temporally persistent changes.
For each spatial position $(i, j)$, we extract a temporal feature vector:

\begin{equation}
    \mathbf{v}_{i,j} = [\mathbf{F}_1[i,j]; \mathbf{F}_2[i,j]; \ldots; \mathbf{F}_T[i,j]] \in \mathbb{R}^{T \times d}
\end{equation}

A transformer encoder with $L$ layers processes this temporal sequence:

\begin{equation}
    \hat{\mathbf{v}}_{i,j} = \text{TransformerEnc}(\mathbf{v}_{i,j} + \mathbf{PE})
\end{equation}

\noindent where $\mathbf{PE} \in \mathbb{R}^{T \times d}$ encodes absolute temporal position (day of year, allowing the model to learn phenological patterns).

The attention mechanism naturally learns to:
\begin{itemize}[leftmargin=*]
    \item Weight cloud-free observations more heavily than cloud-contaminated ones.
    \item Distinguish seasonal vegetation cycles (high attention to phenological outliers) from persistent change (consistent signal across recent timesteps).
    \item Identify gradual degradation patterns that no single bi-temporal comparison would detect.
\end{itemize}

\subsection{Change Decoder}

The change decoder produces two outputs from the attended temporal features:

\begin{equation}
    \hat{c}_{i,j} = D_{\text{detect}}(\hat{\mathbf{v}}_{i,j}) \in [0, 1] \qquad \hat{y}_{i,j} = D_{\text{classify}}(\hat{\mathbf{v}}_{i,j}) \in \mathbb{R}^{C}
\end{equation}

\noindent where $\hat{c}$ is the change probability and $\hat{y}$ is the change type classification logits over $C$ ecological change categories:

\begin{table}[H]
\centering
\caption{Ecological change categories.}
\label{tab:change_categories}
\begin{tabularx}{\textwidth}{@{}llX@{}}
\toprule
\textbf{Category} & \textbf{Code} & \textbf{Description} \\
\midrule
Deforestation & DEF & Conversion of forest to non-forest \\
Degradation & DEG & Reduction in canopy cover without complete removal \\
Agricultural Expansion & AGR & Conversion to cropland or pasture \\
Urban Expansion & URB & Conversion to built environment \\
Mining & MIN & Surface mining and quarrying \\
Fire & FIR & Wildfire or prescribed burning \\
Flooding & FLD & Inundation events (natural or anthropogenic) \\
Reforestation & REF & Regrowth or active replanting \\
No Change & NOC & Stable land cover \\
\bottomrule
\end{tabularx}
\end{table}

\subsection{Training}

The loss function combines binary change detection and multi-class classification:

\begin{equation}
    \mathcal{L} = \mathcal{L}_{\text{BCE}}(\hat{c}, c) + \lambda \cdot \mathbb{1}[c=1] \cdot \mathcal{L}_{\text{CE}}(\hat{y}, y)
\end{equation}

\noindent where $c$ is the ground-truth change label, $y$ is the change type label, and the classification loss is applied only to changed pixels.
We use focal loss variants of both BCE and CE to handle the extreme class imbalance (most pixels are unchanged).

\subsection{Training Dataset}

We construct a training dataset from Sentinel-2 imagery covering 12 biomes:

\begin{table}[H]
\centering
\caption{Training dataset composition by biome.}
\label{tab:training_data}
\begin{tabular}{@{}lcccc@{}}
\toprule
\textbf{Biome} & \textbf{Tiles} & \textbf{Change Pixels (K)} & \textbf{Annotation Source} & \textbf{Years} \\
\midrule
Tropical Forest & 2,847 & 1,234 & GFC + expert & 2019--2025 \\
Temperate Forest & 1,623 & 487 & GFC + expert & 2019--2025 \\
Boreal Forest & 894 & 312 & GFC + FIRMS & 2019--2025 \\
Savanna & 1,247 & 567 & Expert annotation & 2020--2025 \\
Grassland & 782 & 234 & CLC + expert & 2020--2025 \\
Wetland & 534 & 187 & Expert annotation & 2020--2025 \\
Mangrove & 412 & 156 & GMW + expert & 2020--2025 \\
Coral Reef (coastal) & 287 & 89 & Allen Coral Atlas & 2020--2025 \\
Tundra & 198 & 67 & Expert annotation & 2021--2025 \\
Desert Margin & 345 & 123 & Expert annotation & 2020--2025 \\
Montane & 467 & 178 & Expert annotation & 2020--2025 \\
Agricultural Frontier & 1,834 & 892 & MapBiomas + expert & 2019--2025 \\
\midrule
\textbf{Total} & 11,470 & 4,526 & & \\
\bottomrule
\end{tabular}
\end{table}

% ============================================================
% 4. Habitat Fragmentation Index
% ============================================================
